\subsection{Classes from retained source functions}

Let $c$ be an already guarded scalar source function that is rational
at the selected rational character, and put
\begin{equation}\label{eq:adaptiveclass}
 q=2+wc^2,\qquad
 \delta_c(V)=(q^2-2q-1)^2+4q(q^2-1)V.
\end{equation}
The function may use the existing slope or point coordinates, including
finite point-algebra generators whose selected character is rational.
A raw $K$-valued function is not permitted in place of this rational
scalar. No new unrestricted coordinate is introduced by this definition.

\begin{proposition}\label{prop:adaptiveclass}
The class in \eqref{eq:adaptiveclass} preserves the pole implication
and the coordinate census of Theorem~\ref{thm:census}. It is strictly
positive at every prepared character. At every selected contact prime
$p>7$ with positive odd $v_p(w)$ and $V\equiv1$ or $2\pmod p$,
it is a nonzero local square, for every rational value of $c$.
These statements require no bound on the height or denominator of $c$.
\end{proposition}
\begin{proof}
At a hypothetical bad pole, write $\nu=v_p(w)>0$ odd and
$m=v_p(V)>0$. If $c=0$, then $q=2$ and $\delta_c=1+24V$.
Otherwise $v_p(wc^2)$ is odd and nonzero. If it is positive,
$q\equiv2\pmod p$, so $\delta_c\equiv1\pmod p$.
If it is negative, put $e=v_p(q)<0$. The two summands defining
$\delta_c$ have valuations $4e$ and $3e+m$, respectively.
Since $3e+m>4e$, the first, a nonzero square, dominates.
The class is therefore a local square, in particular of even
valuation, in both cases.

At a selected contact $v_p(V)=0$. The same negative-$e$ argument
applies, while the other case reduces to $1+24V$, with residues
$25$ and $49$. They are nonzero squares for $p>7$.

At a prepared real character, $w>0$, $V>0$ and $q\ge2$.
Consequently $4q(q^2-1)V>0$, giving strict positivity.
All new denominators are already denominators of $w,V,c$.
Proposition~\ref{prop:class} and the existing encoder now apply.
\end{proof}

The reference and contact values have the identities
\[
 \delta_c(0)=(q^2-2q-1)^2,\quad
 \delta_c(1)=(q^2+1)^2,\quad
 \delta_c(2)=(q^2+2q-1)^2.
\]
In particular the explicit choice $c=\lambda$, the retained common
slope, is compatible with fixed-place preparation. In a conditional
block with $t=a\lambda$, letting $\lambda$ tend to infinity at a
fixed place gives
\[
 w\lambda^2\longrightarrow\frac1{4a^2},\qquad
 q\longrightarrow2+\frac1{4a^2},\qquad V\longrightarrow2.
\]
The limiting value of $q$ is rational, and $q^2+2q-1$ has no rational
zero. Thus the limiting class is a nonzero rational square; the local
square neighbourhood is open, including at two. In the general source,
the prepared weight tends to a fixed nonzero value as the slope
becomes large at the fixed place. Then $v_p(q)$ tends to minus infinity
and
\[
 \frac{\delta_c(V)}{(q^2-2q-1)^2}
 =1+\frac{4q(q^2-1)V}{(q^2-2q-1)^2}\longrightarrow1.
\]
The same square-neighbourhood argument applies while retaining the
already prepared trace-conic square condition. The source's arbitrary
fixed denominator precision permits this refinement in its original
order of choices.

Any fixed finite selection of the functions in
\eqref{eq:adaptiveclass}, together with the earlier classes, may be
compiled by the same finite product of branch encoders. This allows
the square class to depend on the actual slope and on auxiliary
points without a lost-contact or height penalty. It still requires
simultaneous global rational realization on actual detector tuples.
